\chapter{Why APFS Resists MFT-Style Indexing}

WizTree owes its speed to a peculiarity of NTFS: the Master File Table.
The MFT is a single, flat, contiguous structure in which each file is
one record. Indexing an NTFS volume reduces to reading that table
sequentially. Subdirectory recursion, repeated stat calls, and
path-by-path traversal are all unnecessary.

APFS provides nothing of the kind, by design.

\section{What APFS Provides Instead}

An APFS container is a copy-on-write region of transactional B-trees.
There is no flat metadata table to read; there are at least three
trees the indexer must touch for every file, in a specific order, all
under a single chosen transaction state.

\textbf{Spec.} A container holds:

\begin{itemize}
  \item a stable block-zero locator that points at the
    \emph{checkpoint descriptor area},
  \item the descriptor area itself, which carries a ring of container
    superblocks --- one per transaction --- each of which is a
    self-consistent view of the container at that transaction,
  \item the \emph{container object map} (OMAP), a B-tree that resolves
    virtual object identifiers to physical block addresses under a
    chosen transaction,
  \item one or more \emph{volume superblocks}, reached through the
    container OMAP,
  \item per-volume OMAPs that resolve the volume's virtual objects,
  \item per-volume \emph{file-system trees}, B-trees of variable-length
    records keyed by \texttt{(object\_id, type)} that hold directory
    entries, inodes, extended attributes, hard-link records, and data
    stream metadata.
\end{itemize}

Every metadata read is therefore at least three indirections:
container, volume, FS tree. The lookups themselves are cheap, but each
of them needs the correct OMAP and the correct transaction context.
The naive question --- \emph{where is the MFT?} --- has no answer.

\section{The Right Question}

The right question is structural:

\begin{quote}
What is the minimum set of structures we must traverse, in what
order, under what state pin, to correctly enumerate one volume?
\end{quote}

The answer fills most of this manual. In compressed form:

\begin{enumerate}
  \item Locate the descriptor area from block zero.
  \item Select one checkpoint from that area --- the highest-XID
    checksum-valid container superblock.
  \item Walk the checkpoint map; resolve the container OMAP.
  \item Look up the target volume superblock through the container
    OMAP.
  \item Open the volume OMAP; look up the FS-tree root.
  \item Walk the FS tree, resolving internal-node child references
    through the volume OMAP at every level.
  \item Decode the FS-tree leaf records into directory entries, inodes,
    xattrs, sibling links, and dstreams.
  \item Reconstruct paths, file identity, and logical size from those
    records.
\end{enumerate}

Each step is independently load-bearing. Skip the OMAP lookups at the
FS-tree internal level (chapter~5) and you read random unallocated
blocks. Choose the wrong checkpoint (chapter~3) and you read an
incoherent transaction. Mis-decode an xfield blob (chapter~6) and you
silently emit wrong sizes.

\section{A Narrow Supported Core}

The manual's scope is deliberately narrow:

\begin{itemize}
  \item one APFS volume,
  \item at one pinned transaction state,
  \item correct namespace,
  \item per-file logical size, with per-file allocated size where an
    on-disk field has been oracle-validated,
  \item fail closed outside the tested allowlist.
\end{itemize}

The scope is narrow because the cost of being approximately right is
unbounded: aggregates derived from silently wrong rows propagate into
treemaps, into ``where did my disk space go'' decisions, and into
filesystem operations the user takes based on the result. The cost of
being narrowly right and failing closed everywhere else is bounded: the
fallback path produces a correct answer through supported APIs, more
slowly.

This is not a hedge. It is the way to read APFS metadata at all without
silently corrupting downstream totals.

\section{Why Live Latest Cannot Be The Default}

A particular instinct deserves to be killed early. The instinct says:
"just read the latest checkpoint each time; the volume will be
self-consistent enough." It is wrong on mounted volumes that the user
is using.

\textbf{Observation.} On a mounted lab image, writes from a concurrent
process advance the visible latest checkpoint underneath the
reader's eyes. Reading "latest" twice may return two different
transactions; mixing objects across them produces an incoherent
namespace. The result has no oracle to compare against because no
oracle ever saw the mix. The probe that demonstrated this
(\texttt{EX-05}) is the load-bearing counterexample for the design.

Live latest is therefore not a supportable mode. The supported raw mode
requires a state that can be pinned and held: a detached image, an
explicitly stable APFS source, or a tightly controlled lab volume.
Outside that allowlist the indexer falls back to a supported
operating-system API.

\section{What This Manual Does Not Promise}

Several APFS surfaces are partially or wholly outside the manual's
scope, not because they are unimportant but because they have not yet
been validated end to end. They appear later, in chapter~13, as open
questions:

\begin{itemize}
  \item exclusive and shared bytes per inode, snapshot-retained bytes,
    and any size column that requires the extent-reference tree to be
    joined with per-inode extents;
  \item persistent incremental caches that reuse prior subtree
    summaries (chapter~10 develops the identity rules that any such
    cache would need);
  \item the Finder-visible ``/'' namespace assembled from system and
    data volumes, firmlinks, and snapshots on a modern macOS startup
    disk (chapter~11 documents the conservative firmlink-overlay
    dedup the fallback walker performs; the merged Finder semantics
    proper remain unclaimed).
\end{itemize}

A claim from this manual extends only to the things the manual
explicitly proves.
